\subsection{搜索定位策略}

为每个频道维护状态
\[
  z_c\in\{\text{未知},\text{已发现},\text{待复测},\text{已定位},\text{已清除}\},
  \qquad c=1,\ldots,20.
\]
机器狗按照覆盖目标圆的航点序列移动。在航点处扫描尚未清除的频道；一旦发现信号，记录检测点与示向度，并依据问题二策略选取复测点。两次或多次测向形成的定位区域足够小时，前往其中心或最小覆盖圆圆心。若进入 \(20\,\mathrm{m}\) 光学范围则执行精确定位和清除；若定位失败，则在定位区域边界附近补测并更新约束。

建议按以下优先级选择动作：
\begin{enumerate}
  \item 可立即光学定位并清除的频道；
  \item 已发现且一次补测即可显著缩小定位区域的频道；
  \item 沿当前覆盖路线顺路可检测的未知频道；
  \item 尚未覆盖的空间航点与频道。
\end{enumerate}

一次动作序列的虚拟时间为
\begin{equation}
  T_{\mathrm{all}}=
  \sum_k\frac{\|P_{k+1}-P_k\|_2}{5}
  +N_{\mathrm{switch}}+5N_{\mathrm{detect}}
  +3N_{\mathrm{optical}}+2N_{\mathrm{clear}}.
  \label{eq:time-p3}
\end{equation}
